\section{Scaled variables and ground-state estimates}\label{set:section}\label{sec:setup}

We work in units in which the particle density is large and the product
of density and scattering length stays in a fixed positive interval.
The momentum scale in the theorem is then of order one. The proof must
identify the distribution of excitations on that scale and control the
particles at longer wavelengths. We introduce the spatial projections
that separate these tasks, specify the order of limits, and state the
ground-state estimates from \cite{H} used in the spatial comparison.

\subsection{The form domain and the constant orbital}

For an integer $K\ge100$, put $\Lambda_K=(\R/K\Z)^3$ and $V=K^3$.
The exclusion distance in these units is denoted by $r$.
We use the allowed configuration set and Dirichlet form from the
introduction, with $(L,a)$ replaced by $(K,r)$, and write $E_n$ for the
lowest $n$-particle energy on this same torus with this same exclusion
distance. Set $E_0=0$. When evaluating a nonsymmetric function, $q$
denotes the same Dirichlet form without the bosonic restriction.
All functions are extended by zero to the full configuration torus.
In particular, a normalized bosonic ground vector
$\Phi$ has $\|\Phi\|_2=1$ on $\Lambda_K^N$.

The scaled parameters are
\begin{equation}\label{set:scales}
 D=\frac NV,\qquad \alpha=Dr,\qquad
 \alpha\in I_\alpha:=[\alpha_0/2,2\alpha_0],\qquad
 Dr^3=\frac{\alpha^3}{D^2}.
\end{equation}
Here $\alpha_0$ is a sufficiently large absolute constant, fixed as in
the spatial estimates of \cite{H}. We subsequently increase the lower
threshold on $D$ as needed. Constants may depend on the fixed interval
$I_\alpha$ and on auxiliary parameters already chosen; they are
independent of $K,N$ and the ground vector.

Write $Y=(x_2,\ldots,x_N)$, let $u_0=V^{-1/2}$ on $\Lambda_K$, and
let $P_0=|u_0\rangle\langle u_0|$. For a pure state we abbreviate the
occupation as
\begin{equation}\label{set:occupation}
 B(\Phi)=\|(P_0)_1\Phi\|_2^2
   =\frac1V\int_{\Lambda_K^{N-1}}
       \left|\int_{\Lambda_K}\Phi(y,Y)\,\dd y\right|^2\dd Y.
\end{equation}
The subscript denotes the coordinate on which a one-particle operator
acts. When $\Phi$ is nonnegative, we also use its configuration law
\begin{equation}\label{set:mu}
 \mu(\dd X)=\Phi(X)^2\,\dd X.
\end{equation}
For general complex vectors, configuration expectations use
$|\Phi|^2\,\dd X$; any restriction to nonnegative vectors will be stated.

\begin{lemma}[Ground spaces and allowed slices]\label{set:ground-space}
For every nonempty allowed configuration set with nontrivial form domain,
the full labelled and bosonic Dirichlet operators have compact resolvent
and the same lowest energy. A bosonic ground vector is therefore also a
ground vector of the full labelled operator. Its modulus is a
nonnegative bosonic ground vector. A form-domain function restricts, for
almost every choice of some fixed coordinates, to the Dirichlet form
domain of the open allowed set in the remaining coordinates. Zero
extension of such a slice remains admissible if obstacles contributed
by the fixed coordinates are removed.
\end{lemma}

\begin{proof}
Zero extension sends $H_0^1$ of the allowed set into $H^1$ of the compact
configuration torus. The compact Sobolev embedding proves compactness
of the form embedding, and hence compact resolvent. For a labelled
form-domain function $f$, its root-mean-square symmetrization
\[
 F(X)=\left(\frac1{n!}\sum_{\pi\in S_n}|f(\pi X)|^2\right)^{1/2}
\]
has the same norm and no larger Dirichlet energy, by the gradient
inequality for a vector norm. To justify the Dirichlet condition, first
apply this construction to smooth compactly supported approximants.
Their symmetrizations have compact support in the allowed set, admit
interior smooth approximation, and are bounded in $H^1$. Weak closedness
of $H_0^1$ gives the assertion for $f$. The full and bosonic variational
infima consequently agree. The modulus contraction gives the statement
about $|\Phi|$ by the same variational principle.

For slicing, take smooth compactly supported approximants converging in
the form norm. Fubini's theorem gives a subsequence converging in the
Sobolev norm of the unfixed variables for almost every fixed tuple.
Each approximant has compact support in its open allowed slice, proving
the stated Dirichlet membership. Removing fixed obstacles enlarges that
open set, so zero extension preserves membership. This argument does
not require regularity or connectedness of its boundary. These are the
ground-space and slicing arguments of \cite[Lemmas 2.1 and 2.2]{H}.
\end{proof}

For later use, a real bounded Lipschitz multiplier $u$ satisfies
\begin{equation}\label{set:multiplier}
 q[u\Phi]-E_N\|u\Phi\|_2^2
       =\int |\Phi|^2|\nabla u|^2
\end{equation}
for every ground vector $\Phi$. Indeed test its full weak eigenfunction
equation with $u^2\Phi$, take real parts, and expand the two gradients.
Multiplication by $u$ preserves the form domain.

\subsection{Partitions and the two limiting operations}

Choose small positive exponents $u,w$ with
\begin{equation}\label{set:exponents}
 0<w<u/10000.
\end{equation}
The energy estimates will impose a further absolute upper bound on $u$.
We use integers $R,B$ such that $B$ is a multiple of $R$ and
\begin{equation}\label{set:lengths}
 R\asymp D^w,\qquad B\asymp D^u,
 \qquad d=1/j\quad(j\in\N).
\end{equation}
The fixed factors in these lengths will be specified below to match
the length convention in the local energy estimates.
We take $K$ to be a multiple of $B$. This entails no restriction on
the physical thermodynamic sequences, as the final change of units will
show. The partitions into $B$-cells, $R$-cells, unit cells, and $d$-cells
are nested. Denote the partition into cells of side $l$ by
$\mathcal C_l$. One choice is to place all grid faces at
$-1/2+l\Z$ in each coordinate. Unit cells are then
$B_v=v+[-1/2,1/2)^3$, indexed by
$\mathcal L=(\Z/K\Z)^3$.

For $l\in\{B,R,1,d\}$ define the orthogonal projection
\begin{equation}\label{set:cell-projection}
 (P_l f)(y)=\sum_{A\in\mathcal C_l}\1_A(y)
                   \frac1{|A|}\int_A f(z)\,\dd z,
 \qquad W_l=1-P_l.
\end{equation}
Thus $\sum_i(W_l)_i$ counts particles outside the space of functions
constant on each $l$-cell. The notation $\mathrm d\Gamma(A)=\sum_i A_i$
will also be used for a one-particle operator $A$. For a cell $A$ write
$n_A(X)=\sum_i\1_A(x_i)$. An average over cells always means the uniform
average $|\mathcal C_l|^{-1}\sum_{A\in\mathcal C_l}$.

We use the following convention for all asymptotic estimates. In the
inner, thermodynamic limit, $K\to\infty$ and the actual scaled
parameters converge to fixed values $D_*\ge D_0$ and
$\alpha_*\in I_\alpha$. The integers $B,R,j$ are fixed during this
limit. For $\theta\in\{u,w\}$ define the fixed factor
\begin{equation}\label{set:scale-factors}
 b_\theta=\alpha_0^{-(1+3\theta)/2},
\end{equation}
and choose
\begin{equation}\label{set:nested-lengths}
 R=\lceil b_wD_*^w\rceil,
 \qquad B=R\left\lceil\frac{b_uD_*^u}{R}\right\rceil.
\end{equation}
As $D_*\to\infty$, these choices satisfy
$R/(b_wD_*^w)\to1$ and $B/(b_uD_*^u)\to1$; the latter uses $w<u$.
Their comparability with the powers of the actual $D$ holds eventually
along each inner sequence. The factors in \eqref{set:scale-factors}
ensure that the auxiliary density of the local energy theorem stays
within a fixed factor of $D$, uniformly in $\alpha\in I_\alpha$.
The exact calculation is given in Section~\ref{en:section}.

A displayed dilute error $o(1)$ denotes a bound whose thermodynamic
limsup tends to zero as $D_*\to\infty$, uniformly in the ground vector
and in $\alpha_*\in I_\alpha$. If auxiliary parameters such as $j$ are
fixed, the error may depend on them. We write $o(D^{-1})$ for
$D^{-1}o(1)$ in this convention. An explicitly written $o_K(1)$
vanishes along the inner limit, uniformly in the ground vector; its
volume threshold may depend on $D_*,\alpha_*$ and the fixed auxiliary
parameters. Estimates with powers of $D$ have this convention whenever
they use a thermodynamic energy upper bound. At fixed $D_*$, any
$o_K(1)$ can be made smaller than a prescribed positive multiple of
$D_*^{-\beta}$ by increasing the volume threshold. After enlarging the
constant, a bound $CD^{-\beta}$ can therefore be written without an
additional volume remainder when it is understood to hold beyond that
threshold. Estimates stated for every large finite volume do not need
this convention. In particular,
\begin{equation}\label{set:limit-order}
 o(D^{-1})+o_K(1)
 \quad\hbox{implies}\quad
 \lim_{D_*\to\infty}D_*
       \limsup_{K\to\infty}\bigl(o(D^{-1})+o_K(1)\bigr)=0.
\end{equation}
No bound on the speed of the inner convergence is required.
The microcell parameter $j$ will ultimately be chosen to make an error
$\epsilon_j$ small. It is chosen before the dilute threshold and the
volume threshold; the exponents $u,w$ and the length $R$ need not change
when $j$ is increased.

\subsection{Spatial inputs from the condensation theorem}

The companion condensation theorem \cite{H} supplies the removal-energy
and spatial-count estimates below. They provide particles with which to
construct the path comparisons and bound the cost of discarding crowded
or poorly populated regions. The count estimates concern the law of a
nonnegative ground function; the energy argument and the final momentum
statement also apply to complex ground vectors.

\begin{proposition}[Removal energies]\label{set:removal}
For all sufficiently large $D$, uniformly in \eqref{set:scales},
\begin{equation}\label{set:removal-bound}
 E_N-E_{N-m}\le C\alpha m\qquad(0\le m\le N).
\end{equation}
At fixed $K,r$, the ratio $E_n/n$ is nondecreasing among positive
particle numbers having nonempty form domain.
\end{proposition}

This is \cite[Corollary 3.4 and Equation (3.9)]{H}. Its insertion
argument normalizes the new orbital separately for each old
configuration; the old kinetic energy therefore incurs no
multiplicative loss. Together with the full ground equation, this
estimate is the input to deletion of selected path interactions.

\begin{proposition}[Joint low-count bound]\label{set:low-counts}
One can fix $\alpha_0$ sufficiently large and then $\ell_*>0$
sufficiently small so that every normalized nonnegative bosonic ground
function satisfies, for every $A\subset\mathcal L$,
\begin{equation}\label{set:low-count-bound}
 \mu\{n_{B_v}<\ell_*D\text{ for every }v\in A\}
           \le (C/D^2)^{|A|}.
\end{equation}
The estimate holds for all sufficiently large $D$, uniformly in $K$
and the ground function.
\end{proposition}

This is \cite[Proposition 3.5 and Equation (3.10)]{H}. The constants in this
proposition are fixed before the subsequent path cutoffs are chosen.

\begin{proposition}[Count exponential moments]\label{set:count-mgf}
For $A\subset\mathcal L$ set
\[
 f_A(X)=\sum_{i=1}^N e^{-d_K(x_i,A)},
 \qquad f_{\varnothing}=0,
\]
where the distance is to the set of site centers. Every normalized
nonnegative bosonic ground function obeys
\begin{equation}\label{set:count-mgf-bound}
 \log\int e^{z f_A}\,\dd\mu\le CzD|A|
               \qquad(0\le z\le1).
\end{equation}
The constants are uniform for $K\ge100$, $D$ sufficiently large, and
$\alpha\in I_\alpha$.
\end{proposition}

This is \cite[Proposition 3.7 and Equation (3.18)]{H}.
For example, since $n_{B_v}\le e^{\sqrt3/2}f_{\{v\}}$, it gives
\begin{equation}\label{set:count-tail}
 \mu\{n_{B_v}>C_1D\}\le e^{-cD}
\end{equation}
once $C_1$ is a sufficiently large constant. It also bounds every fixed
moment of a count in a bounded region by $C_kD^k$. Hard packing gives
$n_{B_v}\le C r^{-3}\le C D^3$. Consequently polynomial factors in
local counts do not affect the negligibility of \eqref{set:count-tail}.
For an $R$-cell, a union bound involves only $O(R^3)$ unit cells. These
observations justify removing bounded-count truncations in the moment
estimates proved in Section~\ref{en:consequences}.

\begin{proposition}[Joint close-neighbor tail]\label{set:close-tail}
For each particle label $i$ let
$d_i(X)=\min_{k\ne i}d_K(x_i,x_k)$, with the value $+\infty$ if $N=1$.
Given $\theta>0$, one can choose $b_*>0$
sufficiently small, and then $D$ sufficiently large, so that the set
\[
 \mathcal Q(X)=\left\{v\in\mathcal L:
    \#\{i:d_K(x_i,v)\le5,\ d_i(X)\le b_*D^{-1/3}\}
          \ge\theta D\right\}
\]
satisfies, for every $A\subset\mathcal L$,
\begin{equation}\label{set:close-tail-bound}
 \mu\{A\subset\mathcal Q\}\le p_D^{|A|},
 \qquad
 p_D=\bigl(C_\theta b_*^{-2}D^{-4/3}\bigr)^{1/q_*}\longrightarrow0,
\end{equation}
where $q_*$ is an absolute positive integer. The uniformity is the same
as in Proposition~\ref{set:count-mgf}.
\end{proposition}

This is \cite[Proposition 3.7 and Equation (3.20)]{H}.
It remains true if the distance cutoff is decreased. Neighborhoods of
any other fixed radius can be covered by finitely many radius-$5$
neighborhoods; reducing $\theta$ and taking a bounded coloring gives
the corresponding joint bound. The constants may then depend on that
fixed radius. All these estimates are under the original ground-state
law. They are not assertions about a law conditioned on later path
certificates.
